\documentclass[12pt]{article}
\usepackage{float}
\usepackage{graphicx} % Required for inserting images
\usepackage{fontspec}
\usepackage{apacite}
\setmainfont[
  Path=/usr/local/share/fonts/msttcorefonts/,
  UprightFont=Times_New_Roman.ttf,
  BoldFont=Times_New_Roman_Bold.ttf,
  ItalicFont=Times_New_Roman_Italic.ttf,
  BoldItalicFont=Times_New_Roman_Bold_Italic.ttf
]{Times New Roman}
\usepackage[letterpaper,top=2cm,bottom=2cm,left=2cm,right=2cm,marginparwidth=1.5cm]{geometry}

\title{Weekly assignment 3: Identifying organizational cybersecurity issues related to cybersecurity culture}
\author{Truong Khang Tobias Nguyen}
\date{}

\begin{document}

\maketitle

As established in the previous two assignments, cybersecurity is not just a technical, but a social question as well. This concept also applies to organizational cybersecurity as issues in that regard can be traced back to an organization's cybersecurity culture.
\newline
In that sense, culture dictates behaviour as even the best technical measures and tools are as effective as the people implementing them, how they act, behave and perceive risks. Identifying cultural issues such as lack of awareness, avoid reporting incidents, poor communication, resistance to change, etc. 
\newline
Since culture is intangible, meaning that quantifying it is difficult, we need to think about different methods to identify such issues.
\newline
One way is to observe the behaviour and actions that actually happen in an organization in regard of their cybersecurity practices and the following methods can used for observation:
\begin{itemize}
\item The confidence of the people in an organization to report incidents without the fear of punishment. If reports are withheld, it may imply a blame-oriented culture.
\item  Identifying the rate of timely updating and patching of systems, especially security updates. Low rates may suggest that security tasks are seen as optional or tiresome instead of critical responsibilities.
\item The amount of unapproved or undocumented applications and services (Shadow IT) can be an indicator whether a culture is accepting of official governance/compliance or not.
\item Track the amount of completion and participation in security training. Low rates suggests a culture that is apathetical to cybersecurity.
\item Simulations like phishing attacks can be used to analyze if the culture posseses a certain level of cybersecurity literacy.
\item Post-incident reviews to reveal whether an incident stems from a technical vulnerability on procedural failures and human errors. Especially human sources can be used to reflect on an organization's cultural maturity.
\item Lastly, one simple method can also be surveys, where an organization measures attributes whether their employees understand the consequences of incidents, if they feel responsible for their data and services and how to improve existing cybersecurity measures.
\end{itemize}
In summary, identifying organizational cybersecurity issues is not just about knowing how to defend assets from a technical standpoint but it's also a human and management challenge. 
\newline
It requires measuring behaviors, actions, and ensuring that the an organization's culture is also reflected in their operations.

\end{document}
